% =============================================================================
% Rapport de synthèse, projet SCR
% Compilation : pdflatex rapport_scr.tex (deux fois)
% Les chiffres et les tableaux proviennent de chiffres.tex, engendré par
% rapport/generer_rapport.py à partir des sorties du modèle.
% =============================================================================
\documentclass[11pt,a4paper]{article}

\usepackage[T1]{fontenc}
\usepackage{textcomp}
\usepackage[utf8]{inputenc}
\IfFileExists{french.ldf}{\usepackage[french]{babel}}{%
  \renewcommand{\tablename}{Tableau}\renewcommand{\figurename}{Figure}
  \renewcommand{\contentsname}{Table des matières}}
\IfFileExists{lmodern.sty}{\usepackage{lmodern}}{\usepackage{mathptmx}}
\usepackage[margin=2.3cm,headheight=14pt]{geometry}
\usepackage{graphicx}
\usepackage{booktabs}
\usepackage{tabularx}
\usepackage{xcolor}
\usepackage{amsmath}
\usepackage{amssymb}
\usepackage{siunitx}
\usepackage{fancyhdr}
\usepackage{titlesec}
\usepackage{caption}
\usepackage{float}
\usepackage{enumitem}
\usepackage{microtype}
\usepackage{mdframed}
\usepackage[hidelinks]{hyperref}

\sisetup{group-separator={\,},group-minimum-digits=4,output-decimal-marker={,},
         detect-weight=true,detect-family=true}
\definecolor{bleufonce}{HTML}{1F3864}
\definecolor{grisclair}{HTML}{F4F6FA}
\definecolor{rouge}{HTML}{C00000}
\definecolor{vertfonce}{HTML}{375623}

\titleformat{\section}{\color{bleufonce}\large\bfseries}{\thesection}{0.7em}{}
\titleformat{\subsection}{\color{bleufonce}\normalsize\bfseries}{\thesubsection}{0.6em}{}
\captionsetup{font=small,labelfont={bf,color=bleufonce}}
\setlist[itemize]{itemsep=2pt,topsep=3pt,leftmargin=1.2em}
\setlist[enumerate]{itemsep=2pt,topsep=3pt,leftmargin=1.4em}

\pagestyle{fancy}
\fancyhf{}
\fancyhead[L]{\small\color{bleufonce}Calcul complet du SCR, compagnie composite}
\fancyhead[R]{\small\color{bleufonce}\thepage}
\renewcommand{\headrulewidth}{0.4pt}

\newmdenv[backgroundcolor=grisclair,linecolor=bleufonce,linewidth=0.8pt,
          innertopmargin=6pt,innerbottommargin=6pt,skipabove=8pt,skipbelow=8pt]{encadre}
\newcommand{\apport}[1]{\par\smallskip\noindent\textcolor{vertfonce}{\textbf{Apport de cette
version.}} #1\par\smallskip}
\newcommand{\Meuro}{\text{M\texteuro}}

\input{chiffres}

% =============================================================================
\begin{document}

\begin{titlepage}
\centering
\vspace*{2.5cm}
{\color{bleufonce}\Huge\bfseries Calcul complet du SCR\par}
\vspace{0.6cm}
{\color{bleufonce}\LARGE d'une compagnie d'assurance composite\par}
\vspace{1.2cm}
{\large Formule standard, modèle ALM stochastique, régime 2027,\\[2pt]
modèle interne partiel et ORSA\par}
\vspace{2.2cm}
{\large\bfseries Frege \textsc{Meli Kouyem}\par}
\vspace{0.3cm}
{Master 2 Actuariat, ISFA, Université Claude Bernard Lyon 1\par}
\vspace{2.5cm}
\begin{minipage}{0.82\linewidth}\small\itshape
Projet personnel. Les calculs portent sur une compagnie fictive, calibrée sur des données
publiques réelles : le rapport sur la solvabilité et la situation financière 2025 du groupe
MACSF, les statistiques de placements de l'ACPR, la courbe des taux publiée par l'EIOPA et les
tables de mortalité réglementaires françaises. Date d'évaluation : 31 décembre 2025.
\end{minipage}
\vfill
{\small Environ \num{3700} lignes de Python, 60 tests automatisés\par}
\end{titlepage}

\tableofcontents
\newpage

% =============================================================================
\section{Ce que contient ce rapport}

Le projet reconstitue, de bout en bout, le capital de solvabilité requis d'une compagnie
composite : de la construction des données jusqu'au ratio projeté à cinq ans. Il est organisé en
cinq versions successives, chacune levant une limite de la précédente.

Chaque version est présentée de la même façon : la question à laquelle elle répond, les méthodes
retenues, les hypothèses structurantes, les analyses clés, puis l'apport par rapport à la version
antérieure. Un lecteur pressé peut se limiter aux encadrés de résultats et aux paragraphes
intitulés « apport de cette version ».

\begin{encadre}
\textbf{Résultats d'ensemble.} Pour un actif de \ActifTotal{} \Meuro{}, le SCR ressort à
\ScrVdeux{} \Meuro{} et le ratio de couverture à \RatioVdeux{}. Le risque de marché en porte
l'essentiel, ce qui est le profil attendu d'un assureur à dominante épargne. Le modèle ALM
stochastique dégage une capacité d'absorption par les provisions techniques de \LacTPVdeux{}
\Meuro{}. Le régime applicable en 2027 porte le SCR à \ScrReforme{} \Meuro{} et ramène le ratio à
\RatioReforme{}. Sous stress combiné, le ratio projeté touche un point bas de \RatioOrsaCrise{},
sans franchir le seuil d'alerte.
\end{encadre}

\begin{table}[H]\centering\small
\caption{Les cinq versions et leur apport principal}
\begin{tabularx}{\linewidth}{lXr}
\toprule
\textbf{Version} & \textbf{Apport principal} & \textbf{SCR (\Meuro{})}\\
\midrule
V1 & formule standard complète, provisions déterministes & \ScrVun{}\\
V2 & scénarios économiques, ALM stochastique, absorption mesurée & \ScrVdeux{}\\
V3 & chiffrage de la réforme applicable en 2027 & \ScrReforme{}\\
V4 & modèle interne partiel non-vie et paramètres propres & \ScrVdeux{}\\
V5 & projection à cinq ans et scénarios adverses & \\
\bottomrule
\end{tabularx}
\end{table}

% =============================================================================
\section{La compagnie et ses données}

\subsection{Pourquoi une compagnie fictive}

Aucune base publique ne fournit le portefeuille détaillé d'un assureur : ni les model points de
passif, ni l'inventaire d'actifs ligne à ligne. La démarche retenue consiste donc à construire une
compagnie fictive, mais calibrée sur des données publiques réelles. Les masses de bilan et la
décomposition du SCR viennent du SFCR 2025 du groupe MACSF, ramenées à un facteur d'échelle
d'environ un quart. L'allocation d'actifs suit les statistiques de placements de l'ACPR après mise
en transparence. La courbe des taux et les facteurs de choc proviennent du fichier EIOPA du
31 décembre 2025. Les tables de mortalité sont les tables réglementaires TH et TF 00-02 pour les
garanties décès, TGH05 et TGF05 pour les rentes, ces dernières étant utilisées en générationnel.

\subsection{Le portefeuille en quelques chiffres}

\begin{table}[H]\centering\small
\caption{Volumétrie des données produites}
\begin{tabularx}{\linewidth}{Xr}
\toprule
\textbf{Élément} & \textbf{Valeur}\\
\midrule
Actif total & \ActifTotal{} \Meuro{}\\
Lignes d'actifs dans l'inventaire & \LignesInventaire{}\\
Model points de passif vie et santé & \ModelPoints{}\\
Provisions mathématiques, épargne euro & \PMeuro{} \Meuro{}\\
Provisions mathématiques, unités de compte & \PMuc{} \Meuro{}\\
Capitaux sous risque, temporaire décès & \CapitauxSousRisque{} \Meuro{}\\
Primes annuelles non-vie & \PrimesNonVie{} \Meuro{}\\
Primes annuelles santé, frais de soins & \PrimesSante{} \Meuro{}\\
Triangles de règlements & 5 lignes d'activité sur 10 ans\\
\bottomrule
\end{tabularx}
\end{table}

\begin{figure}[H]\centering
\includegraphics[width=0.84\linewidth]{figures/allocation.png}
\caption{Allocation du portefeuille du fonds euro, après mise en transparence des OPC.}
\end{figure}

\subsection{Trois choix structurants}

\textbf{Structure juridique.} Le principe de spécialisation interdit en France qu'une même entité
exerce à la fois l'assurance vie et l'assurance non-vie. La compagnie est modélisée comme une
entité composite unique, hypothèse explicitement documentée. La scinder en deux entités solo et
calculer un SCR groupe constituerait une extension naturelle.

\textbf{Traçabilité des hypothèses.} Chaque paramètre du fichier de configuration porte une
étiquette indiquant sa nature : donnée réelle issue d'un SFCR, ordre de grandeur de marché, valeur
vérifiée dans le règlement, valeur restant à confronter au texte, ou choix de modélisation. Cette
discipline permet de savoir à tout moment ce qui est solide et ce qui reste provisoire.

\textbf{Reproductibilité.} La graine aléatoire est fixée et propre à chaque module, de sorte
qu'une modification du passif ne déplace pas les tirages de l'actif. Soixante tests automatisés
vérifient les invariants : valorisation exacte des swaps par la courbe, tests de martingale du
générateur de scénarios, égalité comptable à chaque pas de projection, cohérence du SCR avec ses
composantes.

% =============================================================================
\section{V1, la formule standard}

\subsection{La question posée}

Quel capital la formule standard exige-t-elle de cette compagnie, tous modules compris, et la
décomposition obtenue ressemble-t-elle à celle d'un assureur réel de profil comparable ?

\subsection{Méthodes retenues}

\begin{itemize}
\item \textbf{Provisions vie.} Projection model point par model point sur 60 ans :
      revalorisation, rachats structurels, décès, frais, puis actualisation sur la courbe EIOPA.
\item \textbf{Provisions non-vie.} Chain Ladder sur les triangles de règlements, puis
      actualisation des réserves selon la cadence estimée et majoration des frais de gestion.
\item \textbf{Modules de risque.} Chaque choc est évalué en recalculant intégralement les fonds
      propres de base. Le SCR d'un sous-module est la perte de fonds propres qu'il provoque, et
      non la seule perte d'actif.
\item \textbf{Agrégation.} Matrices de corrélation de l'annexe IV, avec le paramètre dépendant du
      scénario de taux retenu.
\end{itemize}

\subsection{Hypothèses fondamentales}

\begin{table}[H]\centering\small
\caption{Hypothèses structurantes de la V1}
\begin{tabularx}{\linewidth}{lX}
\toprule
\textbf{Hypothèse} & \textbf{Valeur et justification}\\
\midrule
Courbe des taux & EIOPA au 31/12/2025, sans correction pour volatilité\\
Revalorisation de l'épargne & maximum entre le taux garanti et 85 \% du rendement attendu\\
Rachats & structurels uniquement, avec le pic fiscal à huit ans\\
Ajustement symétrique & \AjustementSymetrique{} \%, valeur EIOPA du 31/12/2025\\
Absorption par les provisions & implicite, via la règle de revalorisation\\
Absorption fiscale & plafonnée aux impôts différés passifs, majorés de la moitié de l'écart au
plafond réglementaire\\
\bottomrule
\end{tabularx}
\end{table}

\subsection{Analyses clés}

\begin{table}[H]\centering\small
\caption{Best estimate par segment, en \Meuro{}, projection déterministe}
\begin{tabularx}{0.68\linewidth}{Xr}
\toprule
\textbf{Segment} & \textbf{BE}\\
\midrule
\tableauBE
\bottomrule
\end{tabularx}
\end{table}

Le SCR s'établit à \ScrVun{} \Meuro{} pour des fonds propres de base de \FondsPropresVun{}
\Meuro{}, soit un ratio de \RatioVun{}.

\textbf{La décomposition se compare bien au marché.} Exprimée en pourcentage du SCR, elle donne un
risque de marché de \qty{98}{\percent} contre \qty{93}{\percent} pour le groupe de référence, une
souscription vie de \qty{38}{\percent} contre \qty{46}{\percent}, un risque opérationnel de
\qty{4}{\percent} contre \qty{7}{\percent}. Les écarts s'expliquent par une activité non-vie
proportionnellement plus petite et des impôts différés passifs plus faibles. Ce rapprochement
constitue la première validation du modèle.

\textbf{Le Chain Ladder sous-estime les branches longues.} Sur des triangles de dix ans, il manque
de \qty{26}{\percent} à \qty{30}{\percent} des réserves en responsabilité civile, faute d'observer
le facteur de queue. Comme la simulation connaît les ultimes réels, cet écart est mesurable, ce
qui serait impossible sur données réelles.

\textbf{Le rachat massif domine le module vie.} Il n'est appliqué qu'aux contrats dont le rachat
génère effectivement une perte, contrat par contrat, comme le demande l'article 142. Une
application globale sous-estimerait le choc, les contrats favorables venant compenser.

\apport{La V1 pose l'ossature : données, provisions, modules, agrégation, MCR. Elle fournit un
point de comparaison chiffré à toutes les versions suivantes, et met au jour deux limites qui
justifient la suite : l'absorption par la participation aux bénéfices reste implicite, et les
options et garanties financières ne sont pas valorisées.}

% =============================================================================
\section{V2, le modèle ALM stochastique}

\subsection{La question posée}

Combien vaut réellement la capacité d'absorption par les provisions techniques, et quelle est la
valeur temps des options et garanties du fonds euro ? Ces deux montants sont invisibles dans un
calcul déterministe.

\subsection{Méthodes retenues}

\begin{itemize}
\item \textbf{Générateur de scénarios économiques.} Taux court suivant un modèle de Hull-White à
      un facteur, ajusté pour reproduire exactement la courbe initiale. Action et immobilier en
      mouvements browniens géométriques corrélés au taux, sous probabilité risque-neutre.
      Variables antithétiques pour réduire la variance.
\item \textbf{Projection actif-passif.} Le canton épargne euro est projeté par cohortes :
      rendement comptable du portefeuille obligataire avec renouvellement progressif, réalisation
      de plus-values latentes, taux cible indexé sur le dix ans, participation aux bénéfices,
      provision pour participation aux excédents avec reprise sur huit ans, rachats structurels et
      conjoncturels selon les lois de l'ACPR.
\item \textbf{Trois jeux d'actions du management.} Revalorisation dynamique pour le calcul net,
      revalorisation figée au taux cible pour le calcul brut, taux garanti seul pour isoler les
      prestations discrétionnaires futures.
\end{itemize}

\subsection{Hypothèses fondamentales}

\begin{table}[H]\centering\small
\caption{Hypothèses structurantes de la V2}
\begin{tabularx}{\linewidth}{lX}
\toprule
\textbf{Hypothèse} & \textbf{Valeur et justification}\\
\midrule
Nombre de scénarios & \NbScenarios{}, horizon de projection \HorizonAlm{} ans\\
Volatilité du taux court & \qty{1}{\percent}, ordre de grandeur non calibré sur swaptions\\
Volatilité action & \qty{20}{\percent}, ordre de grandeur non calibré sur options\\
Taux cible servi & 80 \% du taux dix ans lissé, majoré de 20 points de base\\
Rachats conjoncturels & lois de l'ACPR, courbe moyenne entre plancher et plafond\\
Allocation du canton euro & proportionnelle aux engagements du fonds, PPE comprise\\
\bottomrule
\end{tabularx}
\end{table}

\subsection{Analyses clés}

\textbf{Le générateur est market-consistent.} Les tests de martingale montrent un écart maximal de
\EcartMartingaleMax{} points de base sur les déflateurs, ce qui valide la simulation avant toute
interprétation des provisions.

\textbf{L'absorption devient mesurable.} Sur les chocs d'actif, la baisse du rendement financier
réduit le taux servi, donc une partie de la perte est reportée sur les assurés. L'ajustement
ressort à \LacTPVdeux{} \Meuro{}, bien en deçà du plafond des prestations discrétionnaires
futures, évaluées à \FDB{} \Meuro{}.

\begin{figure}[H]\centering
\includegraphics[width=0.80\linewidth]{figures/chocs_euro.png}
\caption{Variation du best estimate du fonds euro sous chaque choc, en calcul net et brut
d'absorption. L'écart entre les deux barres est la capacité d'absorption.}
\end{figure}

\textbf{La valeur temps des options est négative.} Le best estimate stochastique s'établit à
\BEeuroVdeux{} \Meuro{}, contre \BEeuroCentral{} \Meuro{} en scénario central : la valeur temps
vaut donc \TVOG{} \Meuro{}. Ce résultat n'est pas une anomalie. La règle de revalorisation, qui
prend le minimum entre le rendement et le taux cible, est concave : l'assureur conserve l'excédent
au-delà de la cible, tandis que le plancher garanti, presque partout nul dans ce portefeuille,
n'est quasiment jamais atteint.

En décalant les taux minimums garantis, on retrouve la forme en cloche caractéristique d'une
valeur temps : maximale lorsque la garantie est à la monnaie, autour de deux points, puis
décroissante lorsqu'elle devient intrinsèque et exerçable dans tous les scénarios.

\begin{figure}[H]\centering
\includegraphics[width=0.72\linewidth]{figures/tvog.png}
\caption{Valeur temps des options et garanties selon le niveau des taux minimums garantis.}
\end{figure}

\textbf{Deux effets se compensent au niveau du ratio.} Le best estimate augmente, donc les fonds
propres passent de \FondsPropresVun{} à \FondsPropresVdeux{} \Meuro{}. Mais le SCR baisse de
\ScrVun{} à \ScrVdeux{} \Meuro{} grâce à l'absorption. Le ratio passe ainsi de \RatioVun{} à
\RatioVdeux{}.

\begin{table}[H]\centering\small
\caption{Tableau S.25.01, capital de solvabilité requis, en \Meuro{}}
\begin{tabularx}{0.78\linewidth}{Xr}
\toprule
\textbf{Poste} & \textbf{Montant}\\
\midrule
\tableauSXXV
\bottomrule
\end{tabularx}
\end{table}

\begin{figure}[H]\centering
\includegraphics[width=0.86\linewidth]{figures/modules.png}
\caption{Décomposition du capital de solvabilité requis. Les modules sont présentés bruts
d'absorption, les ajustements figurant séparément.}
\end{figure}

\apport{La V2 transforme une absorption implicite en un montant chiffré et défendable, et fournit
la valeur temps des options, que la V1 ne pouvait pas voir. Elle installe aussi l'outil dont
toutes les versions suivantes se servent : un générateur de scénarios validé et un modèle ALM
capable de rejouer n'importe quel choc en net et en brut.}

% =============================================================================
\section{V3, l'impact du régime 2027}

\subsection{La question posée}

De combien le ratio bouge-t-il sous le Règlement délégué (UE) 2026/269, adopté le 29 octobre 2025,
publié au Journal officiel le 18 février 2026 et applicable au 30 janvier 2027 ? Et quelles
modifications comptent réellement ?

\subsection{Méthodes retenues}

Le calcul de la V2 est rejoué en empilant les modifications une à une, ce qui produit un pont du
ratio. Chaque étape isole un changement : extrapolation alternative de la courbe, chocs de taux
révisés, corrélation entre taux et spread, corridor de l'ajustement symétrique, facteur de
réassurance non proportionnelle, puis marge de risque. Une marge de risque a également été
recalculée sous le régime actuel, car la V1 la figeait par hypothèse.

\subsection{Hypothèses fondamentales}

\begin{table}[H]\centering\small
\caption{Hypothèses structurantes de la V3}
\begin{tabularx}{\linewidth}{lX}
\toprule
\textbf{Hypothèse} & \textbf{Valeur et justification}\\
\midrule
Extrapolation & premier point de lissage à 20 ans, convergence vers l'UFR à \qty{10}{\percent}\\
Chocs de taux & composante multiplicative et décalage parallèle, tables saisies de mémoire et
signalées comme restant à vérifier\\
Marge de risque & coût du capital ramené de \qty{6}{\percent} à \qty{4.75}{\percent}, avec
dégressivité temporelle\\
Assiette de la marge de risque & SCR des risques non couvrables, donc hors risque de marché\\
\bottomrule
\end{tabularx}
\end{table}

\subsection{Analyses clés}

\begin{table}[H]\centering\small
\caption{Pont du ratio de solvabilité, du régime actuel au régime 2027}
\begin{tabularx}{\linewidth}{Xrrr}
\toprule
\textbf{Étape} & \textbf{SCR} & \textbf{$\Delta$ SCR} & \textbf{Ratio}\\
\midrule
\tableauPont
\bottomrule
\end{tabularx}
\end{table}

\begin{figure}[H]\centering
\includegraphics[width=0.90\linewidth]{figures/pont_2027.png}
\caption{Trajectoire du ratio au fil des modifications réglementaires.}
\end{figure}

\textbf{Les chocs de taux portent tout l'impact défavorable}, soit \ImpactChocsTaux{} \Meuro{}.
Le passage à un choc combinant composante multiplicative et décalage parallèle, avec suppression du
plancher d'un point à la hausse, alourdit nettement la charge. Une analyse de sensibilité montre
que sans le décalage parallèle, l'impact devient nul : tout dépend donc de tables encore non
vérifiées. C'est le point à fiabiliser en priorité.

\textbf{Deux allègements restent sans effet ici.} La corrélation entre taux et spread ramenée de
\num{0.50} à \num{0.25} ne joue qu'en scénario baissier, or la compagnie retient le scénario
haussier. Le corridor élargi de l'ajustement symétrique ne mordrait qu'en cas de marché très
haussier ou très baissier, alors que la valeur du 31 décembre 2025 se situe à l'intérieur des deux
corridors.

\textbf{L'extrapolation alternative est neutre} sur la courbe de fin 2025. Conçue dans
l'environnement de taux bas, elle ne déplace plus la courbe que de quelques points de base au-delà
de trente ans.

\textbf{La marge de risque est le vrai gain}, passant de \MargeRisqueActuelle{} à
\MargeRisqueReforme{} \Meuro{}, soit une baisse de \qty{40}{\percent}.

\apport{La V3 transforme le modèle en outil d'étude d'impact réglementaire. Elle corrige au passage
une faiblesse de la V1, la marge de risque figée par hypothèse, et montre que la lecture habituelle
de la réforme mérite d'être nuancée : ce n'est pas l'extrapolation qui compte pour cette
compagnie, mais les chocs de taux et la marge de risque.}

% =============================================================================
\section{V4, le modèle interne partiel non-vie}

\subsection{La question posée}

La formule standard décrit-elle correctement le risque non-vie de ce portefeuille ? Et que change
une modélisation interne, notamment sur la dépendance entre risques ?

\subsection{Méthodes retenues}

\begin{itemize}
\item \textbf{Réserves.} Bootstrap ODP sur les triangles : rééchantillonnage des résidus de
      Pearson, puis bruit de processus gamma, ce qui donne la distribution complète du boni-mali.
\item \textbf{Primes.} Loi log-normale ajustée sur l'historique des sinistres à primes ultimes
      estimés par Chain Ladder, majorée d'une incertitude d'estimation.
\item \textbf{Catastrophe.} Fréquence de Poisson et sévérité de Pareto généralisée au-delà d'un
      seuil, avec application du traité en excédent de sinistres événement par événement.
\item \textbf{Agrégation.} Copules gaussienne et de Student, ce qui permet de comparer deux
      structures de dépendance à marginales identiques.
\item \textbf{Paramètres propres.} L'écart type spécifique du risque de réserve est déduit du
      bootstrap, puis mélangé au paramètre de marché par le facteur de crédibilité.
\end{itemize}

\subsection{Hypothèses fondamentales}

\begin{table}[H]\centering\small
\caption{Hypothèses structurantes de la V4}
\begin{tabularx}{\linewidth}{lX}
\toprule
\textbf{Hypothèse} & \textbf{Valeur et justification}\\
\midrule
Mesure de risque & valeur en risque à \qty{99.5}{\percent}, en écart à la moyenne\\
Degrés de liberté de la copule de Student & 6, ce qui crée une dépendance de queue marquée\\
Corrélations entre blocs & choix de modélisation, à justifier en annexe d'un dossier réel\\
Sévérité catastrophe & indice de queue de \num{0.35}, paramètre indicatif\\
Crédibilité des paramètres propres & \num{0.74}, correspondant à dix années d'historique\\
\bottomrule
\end{tabularx}
\end{table}

\subsection{Analyses clés}

\begin{figure}[H]\centering
\includegraphics[width=0.90\linewidth]{figures/v4_distribution_non_vie.png}
\caption{Distribution agrégée des charges non-vie à un an, selon la copule retenue.}
\end{figure}

\begin{table}[H]\centering\small
\caption{Comparaison des mesures du risque de souscription non-vie, en \Meuro{}}
\begin{tabularx}{\linewidth}{Xrrr}
\toprule
\textbf{Mesure} & \textbf{SCR non-vie} & \textbf{SCR total} & \textbf{Ratio}\\
\midrule
\tableauModeleInterne
\bottomrule
\end{tabularx}
\end{table}

\textbf{La formule standard est bien calibrée pour ce portefeuille.} Le modèle interne avec copule
de Student, à \ScrNonVieMIS{} \Meuro{}, retrouve presque exactement la formule standard, à
\ScrNonVieFS{} \Meuro{}. Ce n'est pas fortuit : la formule standard a été calibrée pour reproduire
une valeur en risque à \qty{99.5}{\percent} sur un portefeuille moyen, et le nôtre est proche de
la moyenne.

\textbf{La dépendance de queue coûte \qty{6}{\percent} de capital.} Le passage de la copule
gaussienne à celle de Student ajoute \num{5.4} \Meuro{}, alors que les marginales sont identiques.
Cet effet est invisible dans une agrégation par matrice de corrélation, qui ne retient que la
dépendance moyenne.

\textbf{La réassurance agit surtout sur la queue.} Le traité divise la charge catastrophe par
\num{2.6}, de \CatBrute{} \Meuro{} brut à \CatNette{} \Meuro{} net. L'écart entre la valeur en
risque et la perte moyenne au-delà de ce seuil, soit \CatBruteTVaR{} \Meuro{} contre \CatBrute{}
\Meuro{}, illustre l'épaisseur de la queue et l'argument en faveur de la mesure retenue par le
référentiel suisse.

\begin{table}[H]\centering\small
\caption{Écarts types du risque de réserve : propres, standards et retenus}
\begin{tabularx}{0.82\linewidth}{Xrrr}
\toprule
\textbf{Ligne d'activité} & \textbf{Propres} & \textbf{Annexe II} & \textbf{Retenu}\\
\midrule
\tableauUSP
\bottomrule
\end{tabularx}
\end{table}

\textbf{Les paramètres propres doivent être lus avec prudence.} Ils ressortent nettement sous les
paramètres de marché et réduisent le SCR non-vie de \qty{12}{\percent}. Deux réserves s'imposent.
D'une part, les triangles simulés présentent une volatilité inférieure à la moyenne du marché.
D'autre part, et surtout, le bootstrap mesure l'incertitude conditionnellement au modèle Chain
Ladder : il ne capte ni le risque de queue de développement identifié en V1, ni le risque de
modèle.

\textbf{Le rapprochement avec le code CIMA éclaire les référentiels.} La marge de solvabilité
CIMA, calculée sur le même bilan, s'élève à \CimaTotal{} \Meuro{}, soit entre le MCR et le SCR.
La différence de philosophie est nette : la marge CIMA est une exigence forfaitaire indexée sur les
volumes, insensible à l'allocation d'actifs. Or cette compagnie porte l'essentiel de son besoin de
capital en risque de marché, totalement invisible dans la formule CIMA.

\apport{La V4 apporte une mesure de risque construite plutôt que paramétrée, et la possibilité de
tester ce que la formule standard ne voit pas : la dépendance de queue, l'effet fin de la
réassurance, et la spécificité du portefeuille via les paramètres propres. Elle donne aussi un
point de comparaison entre référentiels prudentiels.}

% =============================================================================
\section{V5, l'ORSA}

\subsection{La question posée}

Le ratio tient-il sur l'horizon du plan, et quels risques le mettent réellement sous tension ? La
question n'est plus celle d'un capital à une date, mais celle d'une trajectoire.

\subsection{Méthodes retenues}

Le bilan, le SCR et le ratio sont projetés sur cinq ans à partir du bilan de la V2. Le compte de
résultat est bouclé, de sorte que l'actif se déduit du passif augmenté des fonds propres :
l'égalité comptable est ainsi vraie par construction et vérifiée par un test. Le SCR projeté est
mis à l'échelle d'inducteurs de volume, puis réagrégé chaque année, ce qui est la simplification
usuelle en ORSA. Six scénarios adverses sont testés, complétés par des stress inversés qui
cherchent, par dichotomie, l'amplitude de choc amenant le ratio à la limite d'appétence.

\subsection{Hypothèses fondamentales}

\begin{table}[H]\centering\small
\caption{Hypothèses structurantes de la V5}
\begin{tabularx}{\linewidth}{lX}
\toprule
\textbf{Hypothèse} & \textbf{Valeur et justification}\\
\midrule
Plan d'affaires & décollecte progressive du fonds euro, croissance des primes non-vie de
\qty{3}{\percent}\\
Distribution & \qty{30}{\percent} du résultat net\\
Appétence & limite à \qty{130}{\percent}, alerte à \qty{160}{\percent}, cible à
\qty{200}{\percent}\\
SCR projeté & mis à l'échelle des volumes, sensible à la déformation de l'allocation à hauteur de
\qty{30}{\percent}\\
Actions du management & aucune, ce qui rend les scénarios volontairement sévères\\
\bottomrule
\end{tabularx}
\end{table}

\subsection{Analyses clés}

En scénario central, le ratio s'érode doucement de \RatioVdeux{} à \RatioOrsaCentral{}. Le résultat
net progresse, mais une partie part en dividendes et le SCR croît avec les volumes.

\begin{table}[H]\centering\small
\caption{Scénarios adverses : ratio minimal et statut au regard de l'appétence}
\begin{tabularx}{\linewidth}{Xrrl}
\toprule
\textbf{Scénario} & \textbf{Ratio min.} & \textbf{Ratio final} & \textbf{Statut}\\
\midrule
\tableauOrsa
\bottomrule
\end{tabularx}
\end{table}

\begin{figure}[H]\centering
\includegraphics[width=0.92\linewidth]{figures/v5_trajectoires_orsa.png}
\caption{Trajectoires du ratio de couverture sous les scénarios adverses.}
\end{figure}

\textbf{Le classement des risques est contre-intuitif.} Le krach actions, qui inquiète toujours les
instances de gouvernance, coûte moins que la hausse des taux couplée aux rachats, dont le point bas
s'établit à \RatioOrsaTaux{}, et que la crise combinée, à \RatioOrsaCrise{}. Ce classement recoupe
les conclusions de la V2, où le choc de taux haussier dominait le module marché, et de la V3, où la
réforme l'aggrave encore.

\textbf{Le profil de risque réel dépasse le SCR réglementaire.} Au SCR s'ajoutent les risques mal
captés par la formule standard, à savoir la liquidité, le cyber, le risque de modèle et le climat,
évalués à dire d'expert puis diversifiés. Le besoin global ressort à \BesoinGlobal{} \Meuro{}, soit
\qty{13}{\percent} au-dessus du SCR. C'est le message central de tout exercice ORSA.

\apport{La V5 fait passer le projet du calcul au pilotage. Elle relie les résultats des versions
précédentes à des décisions : seuils d'appétence, capital à injecter, hiérarchie des risques à
surveiller. Elle montre enfin que le risque dominant d'une compagnie n'est pas toujours celui
qu'on redoute le plus.}

% =============================================================================
\section{Ce que chaque version a apporté}

\begin{table}[H]\centering\small
\caption{Progression du projet, version par version}
\begin{tabularx}{\linewidth}{lXX}
\toprule
\textbf{Version} & \textbf{Limite levée} & \textbf{Nouvelle capacité}\\
\midrule
V1 & aucune base de comparaison & calcul complet de la formule standard, validé par rapprochement
avec un SFCR réel\\
V2 & absorption implicite, options non valorisées & capacité d'absorption chiffrée, valeur temps
des options, rachats dynamiques\\
V3 & marge de risque figée, régime unique & étude d'impact réglementaire complète, avec
sensibilités sur les paramètres incertains\\
V4 & risque non-vie réduit à des facteurs & distribution complète des pertes, dépendance de queue,
paramètres propres à l'entreprise\\
V5 & vision à une seule date & trajectoire à cinq ans, scénarios adverses, stress inversés, besoin
global de solvabilité\\
\bottomrule
\end{tabularx}
\end{table}

% =============================================================================
\section{Limites et prolongements}

Toutes les simplifications sont documentées dans le code. Les principales sont les suivantes.

\begin{itemize}
\item \textbf{Paramètres réglementaires.} La courbe, ses facteurs de choc et la table de spread
      proviennent de sources officielles. Restent à confronter au texte les probabilités de défaut,
      les seuils de concentration, les facteurs du MCR et les tables de chocs du régime 2027.
\item \textbf{Calibrage du générateur.} Les volatilités du modèle de taux et de l'action sont des
      ordres de grandeur, et non le résultat d'un calibrage sur prix de swaptions et d'options.
\item \textbf{Catastrophes.} Les modules catastrophe non-vie et santé sont paramétriques. Les
      annexes du règlement, avec leur découpage par zones, restent à implémenter.
\item \textbf{Provisionnement.} L'absence de facteur de queue biaise à la baisse les réserves des
      branches longues, et par conséquent les paramètres propres.
\item \textbf{ORSA.} Le SCR projeté est mis à l'échelle de volumes plutôt que recalculé, et aucune
      action du management n'est déclenchée au franchissement des seuils.
\end{itemize}

Les prolongements naturels sont un modèle interne vie, la scission en entités solo avec calcul d'un
SCR groupe, et une comparaison approfondie avec la norme internationale de l'IAIS ou le référentiel
suisse.

% =============================================================================
\newpage
% =============================================================================
% Annexe A, formulaire : toutes les formules effectivement implémentées.
% Ce fichier est inclus par rapport_scr.tex mais reste lisible seul.
% =============================================================================

\section*{Annexe A, formulaire de la formule standard}
\addcontentsline{toc}{section}{Annexe A, formulaire de la formule standard}

\small
\setlength{\parskip}{4pt}
\setlength{\parindent}{0pt}
Cette annexe recense les formules effectivement implémentées, avec leur référence dans le
règlement délégué (UE) 2015/35. Sauf mention contraire, le SCR d'un sous-module est la perte
de fonds propres de base provoquée par le choc :
\[
\mathrm{SCR}_i \;=\; \max\bigl(0,\; \mathrm{FPB}_0 - \mathrm{FPB}_i\bigr),
\qquad
\mathrm{FPB} \;=\; A - \mathrm{BE} - \mathrm{MR} - \text{autres passifs},
\]
où $A$ est la valeur de marché de l'actif. Les provisions techniques sont intégralement
reprojetées sous chaque choc, ce qui rend le calcul net des actions du management.

% -----------------------------------------------------------------------------
\subsection*{A.1 Provisions techniques et courbe des taux}

\textbf{Best estimate} (art. 28) : $\mathrm{BE} = \sum_{t} \mathrm{FT}_t \, v(t)$, avec
$v(t) = (1+s_t)^{-t}$ le facteur d'actualisation de la courbe sans risque.

\textbf{Extrapolation Smith-Wilson} (régime actuel) : les prix zéro-coupon s'écrivent
$P(t) = e^{-\omega t} + \sum_j \zeta_j\, W(t,u_j)$ avec $\omega = \ln(1+\mathrm{UFR})$ et
\[
W(t,u) = e^{-\omega(t+u)}\Bigl[\alpha\min(t,u) - \tfrac12 e^{-\alpha\max(t,u)}
\bigl(e^{\alpha\min(t,u)} - e^{-\alpha\min(t,u)}\bigr)\Bigr].
\]
Le paramètre $\alpha$ est le plus petit réel supérieur à $0{,}05$ tel que l'écart entre le
forward instantané au point de convergence et $\omega$ n'excède pas un point de base.

\textbf{Extrapolation alternative} (régime 2027) : au-delà du premier point de lissage
$\mathrm{FSP}$, le forward moyen sur $h$ années vaut
\[
f(h) = \mathrm{LLFR} + \bigl(\mathrm{UFR}^{c} - \mathrm{LLFR}\bigr)\,
\frac{1-e^{-a h}}{a h},
\]
où $\mathrm{LLFR}$ est la moyenne pondérée des forwards liquides et $\mathrm{UFR}^{c}$ l'UFR
en composition continue. Sous choc de taux, l'UFR est décalée de $\pm 15$ points de base.

\textbf{Marge de risque} (art. 37 à 39) :
\[
\mathrm{MR} = \mathrm{CoC} \sum_{t\ge 0} \frac{\mathrm{SCR}^{\text{nc}}(t)}{(1+s_{t+1})^{t+1}},
\qquad
\mathrm{SCR}^{\text{nc}}(t) \approx \mathrm{SCR}^{\text{nc}}(0)\,\frac{\mathrm{BE}(t)}{\mathrm{BE}(0)} ,
\]
où $\mathrm{SCR}^{\text{nc}}$ exclut le risque de marché. Le régime 2027 remplace
$\mathrm{CoC}=6\,\%$ par $4{,}75\,\%$ et pondère chaque année par
$\max(\lambda^{t}, \underline{\lambda})$.

% -----------------------------------------------------------------------------
\subsection*{A.2 Risque de marché (art. 164 à 188)}

\textbf{Taux} (art. 166 et 167). Régime actuel : $s_t^{\uparrow} = \max\{s_t(1+f^{\uparrow}_t),\,
s_t + 1\,\%\}$ et $s_t^{\downarrow} = s_t(1-f^{\downarrow}_t)$, sans choc sur les taux négatifs.
Régime 2027 : $s_t^{\uparrow} = s_t(1+b_t) + a_t$ et
$s_t^{\downarrow} = \max\{s_t(1-b_t) - a_t,\, \underline{s}_t\}$, le plancher $\underline{s}_t$
dépendant du terme.

\textbf{Action} (art. 168 à 172) : chocs de $39\,\%$ (type 1) et $49\,\%$ (type 2), majorés de
l'ajustement symétrique $\mathrm{SA}$ ; $22\,\%$ pour les participations stratégiques.
Agrégation :
$\mathrm{SCR}_{\text{act}} = \sqrt{\mathrm{SCR}_1^2 + 2\rho\,\mathrm{SCR}_1\mathrm{SCR}_2
+ \mathrm{SCR}_2^2}$ avec $\rho = 0{,}75$.

\textbf{Immobilier} (art. 174) : choc de $25\,\%$. \quad
\textbf{Change} (art. 188) : $\pm 25\,\%$ par devise, le pire scénario étant retenu.

\textbf{Spread} (art. 176 et 180) : pour une obligation de duration modifiée $d$ et d'échelon
de qualité de crédit $q$,
\[
\mathrm{stress}(q,d) = \min\bigl(a_{q,k} + b_{q,k}\,(d - \underline{d}_k),\; 1\bigr),
\qquad d \ge 1,
\]
$k$ désignant la bande de duration ($\le 5$, $]5,10]$, $]10,15]$, $]15,20]$, $>20$). Les
expositions souveraines de l'EEE libellées en monnaie locale sont exemptées.

\textbf{Concentration} (art. 184 à 187) : pour chaque groupe d'émetteurs $i$,
\[
\mathrm{XS}_i = \max\Bigl(0,\; \frac{E_i}{A^{xl}} - \mathrm{CT}_{q}\Bigr),
\qquad
\mathrm{Conc}_i = A^{xl}\,\mathrm{XS}_i\, g_{q},
\qquad
\mathrm{SCR}_{\text{conc}} = \sqrt{\textstyle\sum_i \mathrm{Conc}_i^2},
\]
$A^{xl}$ étant l'actif hors contrats en unités de compte.

\textbf{Agrégation} (art. 164) :
$\mathrm{SCR}_{\text{mar}} = \sqrt{\mathbf{v}^{\top} \mathbf{C}(A)\, \mathbf{v}}$, où le
paramètre $A$ vaut $0$ si le scénario de taux retenu est la hausse et $0{,}5$ s'il s'agit de la
baisse. Le régime 2027 ramène la corrélation taux-spread à $0{,}25$ en scénario baissier.

% -----------------------------------------------------------------------------
\subsection*{A.3 Défaut de la contrepartie (art. 189 à 202)}

Expositions de type 1 : $\mathrm{LGD}_j = 50\,\% \times (\text{recouvrables} + \text{créances})$
pour la réassurance. La variance de la perte agrégée vaut
\[
V = \sum_{j,k} \frac{\mathrm{PD}_j \mathrm{PD}_k}
{1{,}25(\mathrm{PD}_j+\mathrm{PD}_k) - \mathrm{PD}_j\mathrm{PD}_k}\,
\mathrm{TLGD}_j \mathrm{TLGD}_k
\;+\; \sum_j \frac{1{,}5\,\mathrm{PD}_j(1-\mathrm{PD}_j)}{2{,}5-\mathrm{PD}_j}\,
\sum_{i \in j}\mathrm{LGD}_i^2 ,
\]
puis
\[
\mathrm{SCR}_1 =
\begin{cases}
3\sqrt{V} & \text{si } \sqrt{V} \le 5\,\% \sum \mathrm{LGD},\\
5\sqrt{V} & \text{si } 5\,\% < \sqrt{V} \le 20\,\% \sum \mathrm{LGD},\\
\sum \mathrm{LGD} & \text{sinon.}
\end{cases}
\]
Type 2 : $\mathrm{SCR}_2 = 15\,\%\times$ créances non échues $+\; 90\,\%\times$ créances échues
depuis plus de trois mois. Agrégation :
$\mathrm{SCR}_{\text{déf}} = \sqrt{\mathrm{SCR}_1^2 + 1{,}5\,\mathrm{SCR}_1\mathrm{SCR}_2
+ \mathrm{SCR}_2^2}$.

% -----------------------------------------------------------------------------
\subsection*{A.4 Souscription vie et santé (art. 136 à 163)}

\begin{table}[H]\centering\footnotesize
\begin{tabularx}{\linewidth}{lXl}
\toprule
\textbf{Sous-module} & \textbf{Choc appliqué} & \textbf{Article}\\
\midrule
Mortalité & $q_x \rightarrow 1{,}15\,q_x$, permanent & 137\\
Longévité & $q_x \rightarrow 0{,}80\,q_x$, permanent & 138\\
Invalidité (santé) & taux de sortie d'incapacité réduits de $20\,\%$ & 153\\
Rachat & $\times 1{,}5$ ; $\times 0{,}5$ ; ou rachat massif de $40\,\%$ & 142\\
Frais & $+10\,\%$ de frais et $+1$ point d'inflation & 143\\
Révision & $+3\,\%$ (vie) ou $+4\,\%$ (santé) sur les arrérages & 141, 158\\
Catastrophe & $q_x \rightarrow q_x + 0{,}15$ point la première année & 144\\
\bottomrule
\end{tabularx}
\end{table}

Le choc de rachat massif n'est appliqué qu'aux contrats dont le rachat accroît les provisions,
contrat par contrat. L'agrégation utilise les matrices de l'annexe IV, de forme
$\sqrt{\mathbf{v}^{\top}\mathbf{C}\,\mathbf{v}}$. Le module santé combine ensuite SLT, NSLT et
catastrophe avec les corrélations $\rho_{\text{SLT,NSLT}} = 0{,}5$ et
$\rho_{\cdot,\text{CAT}} = 0{,}25$.

% -----------------------------------------------------------------------------
\subsection*{A.5 Souscription non-vie (art. 114 à 135)}

Pour chaque ligne d'activité $s$, avec $V_{p,s} = \max(P_s, P_{\text{last},s}) + \mathrm{FP}_s$
et $V_{r,s}$ le best estimate de sinistres :
\[
\sigma_s = \frac{\sqrt{(\sigma_{p,s}V_{p,s})^2 + \sigma_{p,s}V_{p,s}\sigma_{r,s}V_{r,s}
+ (\sigma_{r,s}V_{r,s})^2}}{V_{p,s}+V_{r,s}},
\qquad
\sigma = \frac{\sqrt{\sum_{s,t}\mathrm{Corr}_{s,t}\,\sigma_s V_s \sigma_t V_t}}{V},
\]
puis $\mathrm{SCR}_{\text{prim,rés}} = 3\,\sigma\,V$ avec $V = \sum_s V_s$.

\textbf{Paramètres propres à l'entreprise} (art. 218 et annexe XVII) :
$\sigma^{\star} = c\,\sigma^{\mathrm{USP}} + (1-c)\,\sigma^{\text{std}}$, le facteur de
crédibilité $c$ dépendant du nombre d'années d'historique ($c = 0{,}74$ pour dix ans).

\textbf{Catastrophe} : agrégation quadratique des périls naturels et des scénarios d'origine
humaine, puis application du traité en excédent de sinistres,
$L^{\text{net}} = \min(L, \text{priorité}) + \max(0,\, L - \text{priorité} - \text{portée})$.

% -----------------------------------------------------------------------------
\subsection*{A.6 Agrégation finale (art. 87, 103, 204 à 207, 248 à 253)}

\[
\mathrm{BSCR} = \sqrt{\sum_{i,j}\mathrm{Corr}_{i,j}\,\mathrm{SCR}_i\,\mathrm{SCR}_j}
\;+\; \mathrm{SCR}_{\text{incorporels}},
\]
\[
\mathrm{SCR}_{\text{op}} = \min\bigl(30\,\%\times\mathrm{BSCR},\;
\max(\mathrm{Op}_{\text{primes}}, \mathrm{Op}_{\text{prov}})\bigr)
+ 25\,\%\times \mathrm{Dép}_{\text{UC}},
\]
\[
\mathrm{SCR} = \mathrm{BSCR} + \mathrm{SCR}_{\text{op}}
+ \mathrm{Adj}_{\mathrm{TP}} + \mathrm{Adj}_{\mathrm{DT}},
\]
avec $\mathrm{Adj}_{\mathrm{TP}} = -\min\bigl(\mathrm{FDB},\;
\mathrm{BSCR}^{\text{brut}} - \mathrm{BSCR}^{\text{net}}\bigr)$ et
$\mathrm{Adj}_{\mathrm{DT}} = -\min\bigl(\tau\,(\mathrm{BSCR} + \mathrm{SCR}_{\text{op}}
+ \mathrm{Adj}_{\mathrm{TP}}),\; \text{capacité justifiable}\bigr)$, $\tau$ étant le taux
d'impôt. Enfin
\[
\mathrm{MCR} = \max\Bigl\{\min\bigl(\max(\mathrm{MCR}_{\text{lin}},\,25\,\%\,\mathrm{SCR}),\;
45\,\%\,\mathrm{SCR}\bigr),\; \mathrm{AMCR}\Bigr\}.
\]

% -----------------------------------------------------------------------------
\subsection*{A.7 Générateur de scénarios et modèle interne}

\textbf{Hull-White} : $\mathrm{d}r_t = (\theta(t) - a r_t)\mathrm{d}t + \sigma \mathrm{d}W_t$,
de solution exacte à pas annuel
$r_{t+1} = r_t e^{-a} + \alpha(t+1) - \alpha(t)e^{-a}
+ \sigma\sqrt{\tfrac{1-e^{-2a}}{2a}}\,\varepsilon$, avec
$\alpha(t) = f(0,t) + \tfrac{\sigma^2}{2a^2}(1-e^{-at})^2$, et prix zéro-coupon
$P(t,T) = A(t,T)e^{-B(t,T)r_t}$, $B(t,T) = (1-e^{-a(T-t)})/a$. Le contrôle market-consistent
est $\mathbb{E}[D_T] = P(0,T)$ et $\mathbb{E}[D_T S_T]/S_0 = 1$.

\textbf{Valeur temps des options} : $\mathrm{TVOG} = \mathrm{BE}^{\text{stoch}}
- \mathrm{BE}^{\text{central}}$ ; les prestations discrétionnaires futures s'obtiennent par
$\mathrm{FDB} = \mathrm{BE}^{\text{stoch}} - \mathrm{BE}^{\text{garanti}}$.

\textbf{Modèle interne} : $\mathrm{VaR}_{99{,}5\%}(L) - \mathbb{E}[L]$, la perte agrégée $L$
étant construite par copule, $L = \sum_k F_k^{-1}(U_k)$ avec $(U_1,\dots,U_n)$ gaussienne ou de
Student. La TVaR, $\mathbb{E}[L \mid L \ge \mathrm{VaR}]$, est calculée en parallèle pour
mesurer l'épaisseur de la queue.

\setlength{\parskip}{0pt}
\normalsize


\newpage
\section*{Annexe B, organisation du code}
\addcontentsline{toc}{section}{Annexe B, organisation du code}

\begin{table}[H]\centering\small
\begin{tabularx}{\linewidth}{lX}
\toprule
\textbf{Module} & \textbf{Rôle}\\
\midrule
\texttt{scr\_data/} & génération des données : courbe, mortalité, inventaire d'actifs, model
points, triangles, expositions\\
\texttt{scr/passif.py} & best estimate déterministe, Chain Ladder actualisé\\
\texttt{scr/marche.py} & module de marché, courbes choquées, revalorisation ligne à ligne\\
\texttt{scr/souscription.py} & souscription vie et santé\\
\texttt{scr/non\_vie.py} & primes et réserves, rachat, catastrophe\\
\texttt{scr/contrepartie.py} & défaut de type 1 et 2\\
\texttt{scr/agregation.py} & BSCR, opérationnel, ajustements, SCR, MCR, S.25.01\\
\texttt{scr/esg.py} & générateur de scénarios et tests de martingale\\
\texttt{scr/alm.py} & projection actif-passif du fonds euro\\
\texttt{scr/reforme.py} & régime 2027 : extrapolation, chocs, paramètres révisés\\
\texttt{scr/marge\_risque.py} & marge de risque, actuelle et révisée\\
\texttt{scr/modele\_interne.py} & bootstrap, copules, mesures de risque, paramètres propres\\
\texttt{scr/orsa.py} & projection pluriannuelle et besoin global de solvabilité\\
\texttt{tests/} & 60 tests automatisés\\
\bottomrule
\end{tabularx}
\end{table}

\vfill
\begin{center}\small\itshape
Les valeurs et tableaux de ce rapport sont engendrés automatiquement à partir des sorties du
modèle, sans recopie manuelle.
\end{center}

\end{document}
